\chapter{Conclusion}\label{chap:conclusion}

This thesis examines how the transition to globally routable IPv6 addressing affects the exposure of IoT devices in residential networks, and specifically under which router configurations and device settings such a device becomes directly accessible from outside the residential network. To answer this question, a controlled case study was conducted in which a single embedded device, with fixed firmware, service configuration, and addressing behavior, was observed under different perimeter and protocol conditions. This design allowed changes in external reachability to be attributed to the experimental variable rather than to changes in the device itself.
The results show that, in the tested IPv6 setting, external service exposure was determined primarily by the network perimeter rather than by the device. IPv6 restores the architectural possibility of end-to-end connectivity, but whether that possibility becomes exposure depends on the filtering policy enforced at the residential boundary.
The findings also show that accessibility and observability are not equivalent. Even when the firewall discarded all inbound TCP traffic and rendered the HTTP and Telnet services unreachable, the device continued to respond to ICMPv6 and remained identifiable through its EUI-64-derived address. A restrictive perimeter can therefore reduce service accessibility without eliminating discoverability or hardware-derived identification. Once the perimeter permits inbound traffic, however, the behavior of the exposed services becomes directly relevant. The same plaintext administrative access, default credentials, and hidden administrative interface that remained confined to the local network under restrictive perimeter conditions became reachable from outside the residential network without any modification to the device, while the IPv4/NAT baseline kept the same services unreachable without port forwarding.
These observations describe the shift captured in the title of this thesis, from security by obscurity to imminent exposure. The protections that residential IoT deployments have often relied upon, specifically the isolation effect commonly provided by IPv4/NAT and the assumed difficulty of large-scale IPv6 host discovery, were never properties of the devices themselves but external conditions that the transition to IPv6 progressively removes. Within the scope of the conditions examined here, the factors that determined exposure were established primarily by perimeter filtering policies and device-side addressing behavior rather than by actions available to the end user. Reducing this class of risk across the private network ecosystem therefore depends largely on the actors that define those conditions, including ISPs, router manufacturers, and device manufacturers.

\section{Future Work}

Several directions for future research follow from the scope and limitations of this work. The experiment evaluated reachability under known-address conditions, using the global IPv6 address obtained directly from the device. Extending the analysis to the discovery phase would help characterize the complete path from host identification to service access. This includes evaluating whether the same device could be located through adaptive target-generation techniques or inferred through ICMP-related side channels under realistic residential conditions.
A second direction concerns the diversity of residential deployments. The observations reported here are based on a single router model evaluated under two perimeter conditions. Expanding the analysis to additional residential equipment, filtering policies, and deployment models would clarify how representative the observed behavior is across the broader ecosystem. This includes intermediate filtering configurations, ISP-defined defaults, alternative IPv6 transition mechanisms, and access architectures that differ from the native dual-stack environment examined in this work.
The threat model examined here focused on the conditions that enabled external access rather than on the consequences that follow from it. Future work could therefore extend the analysis beyond successful authentication and service access to evaluate post-access activity under controlled conditions, including persistence, lateral movement, and other forms of compromise. Finally, at the device level, additional research is also needed on the adoption of privacy-preserving addressing mechanisms in constrained IoT environments. Evaluating the feasibility, performance impact, and implementation cost of alternatives to EUI-64 in resource-limited devices would help address the identifiability dimension that perimeter controls cannot eliminate.